\documentclass{beamer}
\usepackage{ctex}
\usetheme{SimplePlus}
\setbeamertemplate{footline}{}

\usepackage{amsmath,amssymb,mathtools}
\usepackage{tikz}
\usetikzlibrary{arrows.meta}
\usepackage{tikz-cd}

\usepackage{unicode-math}
%\setmathfont{STIX Two Math}
%\setmathfont{TeX Gyre Termes Math}
%\setmathfont{Libertinus Math}
%\setmathfont{Fira Math}

\usepackage[
  hybrid,                            % 允许在 Markdown 中混用 LaTeX 命令
  texMathDollars                     % 启用 $...$ 和 $$...$$ 数学公式
]{markdown}

\usepackage{graphicx} % Allows including images
\graphicspath{{./images/}}

\title{profunctor}

\date{}

\begin{document}

\begin{frame}
    % Print the title page as the first slide
    \titlepage
\end{frame}

\begin{frame}[fragile]{定义}
设 \(\mathcal C,\mathcal D\) 是范畴。

一个从 \(\mathcal C\) 到 \(\mathcal D\) 的 profunctor，是一个函子
\[
\boxed{
\Phi:\mathcal D^{\mathrm{op}}\times\mathcal C\to\mathbf{Set}
}
\]

记作
\[
\boxed{
\Phi:\mathcal C\nrightarrow\mathcal D
}
\]
\end{frame}

\begin{frame}{Profunctor = “不一定来自同一个范畴的 Hom”}
普通 Hom 是：
$$
\mathcal C(-,-):
\mathcal C^{op}\times\mathcal C\to\mathbf{Set}.
$$

这里两个对象：
$$
c,d
$$

来自同一个范畴。

Profunctor 则把它推广为：
$$
\phi:D^{op}\times C\to\mathbf{Set}.
$$

于是：
$$
\phi(d,c)
$$

就像一个“跨范畴的 Hom”。

所以可以把它理解成：
$$
\boxed{
\text{Profunctor}
\approx
\text{generalized Hom between two categories}.
}
$$

这其实比“二元函子”的定义更有直觉。
\end{frame}

\begin{frame}{Profunctors as categories}
一个 Profunctor 可以直接看成一个“拼接出来的范畴”，其中 \(C\) 和 \(D\) 是两个子范畴，而 Profunctor 给出了它们之间额外的“异质态射”。
\end{frame}

\begin{frame}[shrink]{Functor 是一种特殊的 Profunctor}
如果：
$$
F:C\to D,
$$

Yoneda embedding：
$$
Y_D:D\to\widehat D
$$

把 \(F\) 变成：
$$
Y_D\circ F:C\to\widehat D.
$$

由于：
$$
\widehat D=\mathbf{Set}^{D^{op}},
$$

所以对于 \(c\in C\)，得到一个 presheaf：
$$
Y_D(Fc)=D(-,Fc).
$$

也就是说：
$$
\boxed{
\phi_F(d,c)=D(d,Fc).
}
$$

因此：
$$
\phi_F:D^{op}\times C\to\mathbf{Set}.
$$

这就是一个 Profunctor。材料也明确给出了：
$$
\phi_F=Y_D\circ F.
$$

所以：
$$
\boxed{
F:C\to D
\quad\leadsto\quad
\phi_F(d,c)=D(d,Fc).
}
$$

可以理解成：

\begin{quote} Functor 给每个 \(c\) 指定一个 \(D\) 中的对象 \(Fc\)；Profunctor 则进一步问：任意 \(d\) 到这个 \(Fc\) 有多少态射？\end{quote}
\end{frame}

\begin{frame}[shrink]{Yoneda 和 Profunctor 的联系}
注意：
$$
Y_D(d)=D(-,d).
$$

所以一个对象：
$$
d\in D
$$

实际上可以被表示成一个 presheaf：
$$
D(-,d).
$$

而一个 Profunctor：
$$
\phi:D^{op}\times C\to Set
$$

等价于：
$$
\widehat\phi:C\to\widehat D.
$$

即：
$$
\boxed{
\phi(d,c)
=
\widehat\phi(c)(d).
}
$$

所以你可以把 Profunctor 看成：
$$
\boxed{
C\longrightarrow \widehat D
}
$$

即：

\begin{quote}\textbf{把 \(C\) 中的每一个对象，送到 \(D\) 上的一个 presheaf。}\end{quote}

这实际上是非常重要的视角。
\end{frame}

\begin{frame}[shrink]{Composition of profunctors}
The composite $\psi\phi$ of two profunctors
\[
\phi \colon C \nrightarrow D \quad\text{and}\quad \psi \colon D \nrightarrow E
\]
is given by
\[
\psi\phi = \operatorname{Lan}_{Y_D}(\hat{\psi}) \circ \hat{\phi}
\]
where $\operatorname{Lan}_{Y_D}(\hat{\psi})$ is the left \textcolor{blue}{Kan extension} of the functor $\hat{\psi}$ along the \textcolor{blue}{Yoneda functor} $Y_D \colon D \to \hat{D}$ of $D$ (which to every object $d$ of $D$ associates the functor $D(-,d) \colon D^{\mathrm{op}} \to \mathbf{Set}$).

It can be shown that
\[
(\psi\phi)(e,c) = \Biggl( \coprod_{d\in D} \psi(e,d) \times \phi(d,c) \Biggr) \Big/ \sim
\]
where $\sim$ is the least \textcolor{blue}{equivalence relation} such that $(y',x') \sim (y,x)$ whenever there exists a morphism $v$ in $D$ such that
\[
y' = vy \in \psi(e,d') \quad\text{and}\quad x'v = x \in \phi(d,c).
\]
Equivalently, profunctor composition can be written using a \textcolor{blue}{coend}
\[
(\psi\phi)(e,c) = \int^{d\in D} \psi(e,d) \times \phi(d,c).
\]
\end{frame}

\begin{frame}{Bicategory of profunctors}
Composition of profunctors is associative only up to isomorphism (because the product is not strictly associative in $\mathbf{Set}$). The best one can hope is therefore to build a \textcolor{blue}{bicategory} $\mathbf{Prof}$ whose
\begin{itemize}
  \item 0-cells are \textcolor{blue}{small categories},
  \item 1-cells between two small categories are the profunctors between those categories,
  \item 2-cells between two profunctors are the \textcolor{blue}{natural transformations} between those profunctors.
\end{itemize}
\end{frame}

% 参考文献帧，长列表自动分页
\begin{frame}[allowframebreaks]{参考文献}
\begin{thebibliography}{99}

\bibitem{wiki_profunctor}
\url{https://en.wikipedia.org/wiki/Profunctor}

\end{thebibliography}
\end{frame}
\end{document}